Diffusion-based virtual try-on models produce realistic images of a person wearing a given garment. However, they are mostly evaluated on their final output. Little is known about how that output emerges during denoising. We decode the model's estimate of the clean image at every denoising step and score it on six axes: pose, structure, colour, texture, pattern and stability. We compare CatVTON, OOTDiffusion and IDM-VTON on 48 VITON-HD person-garment pairs, six in each of eight garment categories. All models use the same sampler and are scored over the same garment region at a common resolution. We also decompose every estimate into coarse, mid and fine spatial-frequency bands. All three models fix the pose and coarse structure of the garment within the first few steps and resolve fine detail last. They differ in timing: OOTDiffusion fills in the garment's texture and pattern from the first steps, whereas IDM-VTON and CatVTON only start to do so around step 5 and 10 respectively. Uniform garments reach their final structure within a few steps and then plateau, while repeating patterns keep changing much longer. These differences indicate where in the schedule each model could be improved or combined with another.

